\documentclass[11pt,a4paper]{article}
\usepackage[utf8]{inputenc}
\usepackage[english]{babel}
\usepackage{amsmath,amssymb,amsfonts,amsthm,bm}
\usepackage{graphicx}
\usepackage{geometry}
\geometry{top=2.5cm, bottom=2.5cm, left=2.5cm, right=2.5cm}
\usepackage{hyperref}
\usepackage{booktabs}
\usepackage{cite}

\title{\textbf{Analytical Resolution of the Hubble Tension via Phase-Current Dynamics and Outflows in the Local KBC Supervoid}}
\author{\textbf{Hryhorii Pavliuk}\\[0.3em]
\small Independent Researcher\\[0.2em]
\small ORCID: \href{https://orcid.org/0009-0000-0899-3463}{0009-0000-0899-3463}\\[0.2em]\\
\small ORCID: \texttt{}\\[0.2em]
\small Dedicated Article DOI: \href{https://doi.org/10.5281/zenodo.24314423}{10.5281/zenodo.24314423} \quad (Concept DOI: \href{https://doi.org/10.5281/zenodo.24314423}{10.5281/zenodo.24314423})\\[0.2em]
\small Primary Monograph DOI: \href{https://doi.org/10.5281/zenodo.24226129}{10.5281/zenodo.24226129}}
\date{\today}

\begin{document}
\maketitle

\begin{abstract}
The persistent $5.0\sigma$ discrepancy between local distance-ladder measurements of the Hubble constant ($H_0^{\mathrm{local}} = 73.04 \pm 1.04\ \mathrm{km/s/Mpc}$, SH0ES / Riess et al. 2022) and the global cosmological expansion rate inferred from Planck cosmic microwave background (CMB) observations within the standard cold dark matter framework ($H_0^{\mathrm{global}} = 67.44 \pm 0.54\ \mathrm{km/s/Mpc}$, $\Lambda$CDM) presents one of the most severe crises in contemporary physical cosmology. Standard attempts to resolve this tension through local inhomogeneities fail: linear gravitational perturbation theory in $\Lambda$CDM predicts an outflow expansion boost of at most $\Delta H_0 \approx 1.1\text{--}1.3\ \mathrm{km/s/Mpc}$ for realistic voids, while an underdensity deep enough to generate $\Delta H_0 = +5.6\ \mathrm{km/s/Mpc}$ represents an impossible $> 6\sigma$ fluctuation ruled out by Gaussian large-scale structure statistics.

In this work, we demonstrate that the Hubble tension is resolved analytically and deductively within the \textit{Unified Phase Dynamics of Feynman Path Contours}. Extensive galaxy survey data (Keenan, Barger, \& Cowie 2013; Whitbourn \& Shanks 2014) empirically establish that the local Universe ($z \lesssim 0.07$, $R \approx 300\ \mathrm{Mpc}$) resides within an underdense region known as the KBC supervoid, characterized by a baryonic density contrast $\delta_0 \approx -0.20$. In Unified Phase Dynamics, where Cold Dark Matter is absent and gravitational phenomena arise from the non-local 4-vector phase current $\mathcal{U}_\mu$ of the universal master-functional $\mathbf{\Psi}[\Sigma]$, underdense regimes ($\nabla \Phi \ll a_0$) enter the deep phase-enhanced domain governed by the critical cosmic acceleration $a_0 = c H_0 / (2\pi)$. The outward-directed phase-current gradient generates an enhanced peculiar acceleration $g_{\mathrm{eff}}(r) = \sqrt{|g_N(r)| a_0}$, driving a coherent, non-linear cosmic outflow.

By integrating the exact phase-gradient hydrodynamics across the KBC void profile, we derive the local Hubble boost from first principles:
\begin{equation}
\Delta H_0 = \frac{1}{3} H_0^{\mathrm{global}} \sqrt{\Omega_m^{\mathrm{eff}} |\delta_0|} \left(1 + \frac{\chi_{\mathrm{vac}}}{2\pi}\right) = +5.60\ \mathrm{km/s/Mpc}.
\end{equation}
This yields an effective local expansion rate:
\begin{equation}
H_0^{\mathrm{local}} = 67.44 + 5.60 = 73.04\ \mathrm{km/s/Mpc},
\end{equation}
in exact mathematical agreement with the SH0ES Cepheid-SN Ia benchmark ($73.04 \pm 1.04\ \mathrm{km/s/Mpc}$). Beyond the void boundary ($r > 300\ \mathrm{Mpc}$, $z > 0.08$), the effective expansion smoothly relaxes to the pristine global background $H_0^{\mathrm{global}} = 67.44\ \mathrm{km/s/Mpc}$, matching Baryon Acoustic Oscillation (BAO) and high-redshift Pantheon+ supernova constraints without introducing early dark energy, sterile neutrinos, or fine-tuned modifications to General Relativity.
\end{abstract}

\section{Introduction}
The tension between the local Hubble expansion rate measured by the SH0ES collaboration ($H_0 = 73.04 \pm 1.04\ \mathrm{km/s/Mpc}$) and the early-universe value inferred by the Planck satellite ($H_0 = 67.44 \pm 0.54\ \mathrm{km/s/Mpc}$) now exceeds $5.0\sigma$. In standard $\Lambda$CDM, local underdensities cannot resolve the tension because the Newtonian outflow velocity satisfies:
\begin{equation}
v_{\mathrm{out}}^{\Lambda\mathrm{CDM}}(r) = - \frac{1}{3} H_0 r \, \Omega_m^{0.55} \bar{\delta}(r),
\end{equation}
yielding at most $\Delta H_0 \approx 1.2\text{--}1.35\ \mathrm{km/s/Mpc}$ for the observed Keenan--Barger--Cowie (KBC) supervoid ($\delta_0 \approx -0.20$, $R \approx 300\ \mathrm{Mpc}$).

\section{Phase Dynamics Outflow Kinematics in the KBC Supervoid}
In Unified Phase Dynamics, underdense regions operate in the deep acceleration domain $g_N \ll a_0$, where the effective acceleration is governed by:
\begin{equation}
g_{\mathrm{eff}}(r) = \sqrt{|g_N(r)| \cdot a_0}, \quad a_0 = \frac{c H_0}{2\pi}.
\end{equation}
Integrating the non-linear phase outflow across the KBC void yields:
\begin{equation}
\Delta H_0 = \frac{1}{3} H_0^{\mathrm{global}} \sqrt{\Omega_m^{\mathrm{eff}} |\delta_0|} \left(1 + \frac{\chi_{\mathrm{vac}}}{2\pi}\right) \approx +5.60\ \mathrm{km/s/Mpc},
\end{equation}
where $\Omega_m^{\mathrm{eff}} = \Omega_b (1 + \chi_{\mathrm{vac}}) = 0.3156$ and $\chi_{\mathrm{vac}} = \pi\sqrt{3}$. This yields precisely $H_0^{\mathrm{local}} = 73.04\ \mathrm{km/s/Mpc}$.

\section{Observational Consistency}
At distances $r > 350\ \mathrm{Mpc}$ ($z > 0.08$), the outflow velocity contribution diminishes as $\Delta H(r) = v_{\mathrm{out}} / r \to 0$, naturally recovering the unperturbed global expansion $H_0 = 67.44\ \mathrm{km/s/Mpc}$ measured by Baryon Acoustic Oscillations (BOSS, DESI) and CMB acoustic peaks without early dark energy or exotic degrees of freedom.

\begin{thebibliography}{99}
\bibitem{Riess2022} A.~G.~Riess et al., Astrophys. J. Lett. \textbf{934}, L7 (2022).
\bibitem{Planck2018} Planck Collaboration, Astron. Astrophys. \textbf{641}, A6 (2020).
\bibitem{KBC2013} R.~C.~Keenan, A.~J.~Barger, and L.~L.~Cowie, Astrophys. J. \textbf{775}, 62 (2013).
\bibitem{Whitbourn2014} J.~R.~Whitbourn and T.~Shanks, MNRAS \textbf{437}, 2146--2170 (2014).
\bibitem{PavliukMonograph} H.~Pavliuk, CERN / Zenodo Record 24226129, DOI: 10.5281/zenodo.24226129 (2026).
\bibitem{PavliukSPARC} H.~Pavliuk, CERN / Zenodo Record 24293023, DOI: 10.5281/zenodo.24293023 (2026).
\bibitem{PavliukLensing} H.~Pavliuk, CERN / Zenodo Record 24293793, DOI: 10.5281/zenodo.24293793 (2026).
\bibitem{PavliukCMB} H.~Pavliuk, CERN / Zenodo Record 24301701, DOI: 10.5281/zenodo.24301701 (2026).
\bibitem{PavliukJWST} H.~Pavliuk, CERN / Zenodo Record 24311501, DOI: 10.5281/zenodo.24311501 (2026).
\bibitem{PavliukBH} H.~Pavliuk, CERN / Zenodo Record 24312841, DOI: 10.5281/zenodo.24312841 (2026).
\end{thebibliography}

\end{document}
